\documentclass[a4paper,11pt,dvipdfmx]{jsarticle}

\usepackage[top=25mm,bottom=25mm,left=25mm,right=25mm]{geometry}
\usepackage{amssymb}
\usepackage{listings}
\usepackage{plistings}
\usepackage[dvipdfmx,hidelinks]{hyperref}

\lstset{
  basicstyle=\ttfamily\small,
  frame=single,
  breaklines=true,
  xleftmargin=1em,
  aboveskip=6pt, belowskip=6pt,
  columns=flexible,
}

\newcommand{\cboxon}{$\boxtimes$}
\newcommand{\cboxoff}{$\square$}

\begin{document}

{\LARGE\bfseries 生成AI利用申告書}
\vspace{1em}

\begin{itemize}
  \item 氏名：塩澤 匠生
  \item 学籍番号：25G1065
  \item プログラムファイル名：\texttt{HCK02\_02.ino} / \texttt{HCK02\_02.pde}
  \item 使用したAIツール：ChatGPT (GPT-5)、Claude Code (Opus 4.7)
\end{itemize}

\section*{1. 何に使ったか}
提出課題2 は「前回課題（ハッカソン1 第1回）で作成した楽曲を Processing で
再生しつつ、マイクから入った音の波形を Processing のウィンドウに描画する」
課題である。授業の例題1 では Arduino から Processing へ 1 バイトの値を送り、
円の色を切り替える書き方（\texttt{Serial.write} と \texttt{serialEvent} +
\texttt{read()}）までを学んだが、波形のように「サンプルを連続して送って
時系列に並べて表示する」書き方は扱われていなかった。そこで
\begin{itemize}
  \item Arduino から Processing へマイクのサンプルを連続して渡す方法
  \item Processing 側で受け取ったサンプルを時系列の波形として描画する方法
        （バッファの持ち方、描画ループとの分離）
\end{itemize}
の 2 点を生成AIに相談した。楽曲再生（Minim を使う部分）は前回課題で
書いたコードを流用したため、AIに相談していない。

\section*{2. AIへの指示（プロンプト）}
\begin{lstlisting}
Arduino でマイクの値を analogRead して、Processing のウィンドウに
時系列の波形として描きたいです。
- 授業の例題1 では Serial.write で 1 バイト送って、Processing 側で
  serialEvent から read() を呼ぶ書き方を習いました
- 今回は値を連続で送り続け、ウィンドウの左から右に波形が並ぶように
  描きたいです
Arduino 側はどう送るのが扱いやすいですか? また Processing 側では
受信したサンプルをどんなバッファに貯めて、どう毎フレーム描けば
よいですか?
\end{lstlisting}

\section*{3. AIの出力の使い方}
\begin{itemize}
  \item \cboxoff そのまま使った
  \item \cboxon 一部修正して使った
  \item \cboxoff 参考にしただけ
\end{itemize}
$\rightarrow$ 上で選んだ理由：

AI からは、
\begin{itemize}
  \item Arduino 側は \texttt{Serial.println(d)} で「1 行 1 サンプル」のテキスト
        として送ると、Processing 側で改行区切りで切り出せて扱いやすい
  \item Processing 側は \texttt{port.bufferUntil('\textbackslash n')} と
        \texttt{serialEvent} を組み合わせ、改行が来たときだけ \texttt{readStringUntil}
        で 1 行受け取って整数に変換する
  \item 描画は循環バッファ（\texttt{int[]} と書き込みインデックス）に値を貯め、
        \texttt{draw()} の中で「最新位置から左方向に直近 N サンプル」を読み出して
        線で結ぶ
\end{itemize}
という方針が出てきた。例題1 の \texttt{serialEvent} の使い方とも整合していたので
この骨格はそのまま採用した。一方で、最初に提示されたサンプルは「2 バイトの
バイナリで送る」「描画は \texttt{draw()} の中で \texttt{port.read()} を呼ぶ」
書き方で、デバッグしづらく、また値の取りこぼしが起きそうだったので採用しなかった。

\section*{4. 正しいと確認した方法}
\begin{itemize}
  \item Arduino IDE で \texttt{HCK02\_02.ino} がコンパイルでき、シリアルモニタ
        (921600\,bps) に整数が 1 行ずつ流れることを確認した
  \item Processing IDE で \texttt{HCK02\_02.pde} を起動し、無音時はほぼ平らな線、
        マイクに声を出すと白い線が縦に振れる様子を観察した
  \item キー \texttt{p} を押すと前回課題（ハッカソン1 第1回）と同じ楽曲が
        PC のスピーカから流れ、その音をマイクが拾って波形がより大きく振れる
        ことを確認した
  \item キー \texttt{1}〜\texttt{6} で音色を切り替えてから \texttt{p} を押すと
        異なる音色で再生され、波形の見え方も変わることを確認した
  \item Processing のコードを読み直し、循環バッファのインデックス計算
        \texttt{(writeIdx - DISPLAY\_SAMPLES + samples.length) \% samples.length}
        がなぜ「右端が最新サンプル」になるかを自分で説明できるようにした
\end{itemize}

\section*{5. 問題や間違い}
\begin{itemize}
  \item \cboxon あった
  \item \cboxoff なかった
\end{itemize}

最初に AI が出してきた Processing コードは、\texttt{port.read()} を
\texttt{draw()} の中で呼ぶ書き方だった。これだと値の取りこぼしや
描画ループとのタイミングずれが起きやすそうだったので、授業の例題1 と
同じ \texttt{serialEvent} 方式に揃えた。また、起動直後にゴミ文字が
混ざってきたとき \texttt{Integer.parseInt} が例外を投げる場合があったため、
\texttt{NumberFormatException} を握りつぶす \texttt{try-catch} を入れて
落ちないようにした。

\section*{6. 自分で工夫した点}
\begin{itemize}
  \item Arduino 側 (\texttt{HCK02\_02.ino}) は、提出課題1 のプロッタ用ガイド
        (\texttt{0}, \texttt{1023}) を取り除いて \texttt{Serial.println(d)} 1 本に
        絞り、Processing が一番シンプルに受け取れる形にした
  \item Processing 側 (\texttt{HCK02\_02.pde}) は、前回課題の楽曲再生コード
        (\texttt{HackInstrument} クラス、メロディ・拍・音量の配列、Minim 関連の
        セットアップ) をそのまま組み込み、波形描画と楽曲再生が同じウィンドウ・
        同じキー操作 (\texttt{p}, \texttt{1}--\texttt{6}) で扱えるようにした
  \item 循環バッファの長さをウィンドウ幅 \texttt{width} と揃え、画面幅と
        サンプル数の対応をシンプルにした
  \item 直近の \texttt{DISPLAY\_SAMPLES} 個（60 サンプル）だけを画面いっぱいに
        引き伸ばして描画することで、約 6 周期程度のゆったりした波形になり、
        見やすさを確保した
\end{itemize}

\section*{参考文献}
\begin{itemize}
  \item 有本泰子, 佐波孝彦「ハッカソン1 第2回」講義資料
        (\texttt{HCK1\_02.pdf}, \texttt{HCK1\_02\_ex.pdf})
  \item 生成AI利用申告書テンプレート（ハッカソン1 講義資料）
  \item Arduino Reference: \texttt{Serial.println}, \texttt{analogRead} \\
        \url{https://docs.arduino.cc/language-reference/}
  \item Processing Reference: \texttt{Serial}, \texttt{bufferUntil}, \texttt{serialEvent} \\
        \url{https://processing.org/reference/libraries/serial/}
  \item Minim Reference: \texttt{AudioOutput}, \texttt{Oscil}, \texttt{Line} \\
        \url{https://code.compartmental.net/minim/}
\end{itemize}

\end{document}
